\documentclass{article}
\usepackage{../mathnotezh}

\title{三面角余弦定理}
\author{tbj}
\date{}

\begin{document}

\maketitle

设$OA, OB, OC$为空间中共顶点的三条射线,这样所成的几何结构称为\textbf{三面角}.

考虑这样所成的三个平面角$\angle AOB, \angle BOC, \angle COA$,
以及三个二面角$B \mbox{-} OA \mbox{-} C, C \mbox{-} OB \mbox{-} A, A \mbox{-} OC \mbox{-} B$
(分别记为$\angle OA, \angle OB, \angle OC$).

注意到,将$A, B, C$在其分别所在的射线上滑动,不会改变上述六个角的值.
于是不妨设$OA \perp AB, AC$,则$\angle OA = \angle BAC$.

由$\triangle ABC$中的余弦定理知$\cos \angle BAC = \dfrac{AB^2 + AC^2 - BC^2}{2AB \cdot AC}$.

其中
\begin{align*}
    AB^2 & = OB^2 - OA^2 \\
    AC^2 & = OC^2 - OA^2 \\
    BC^2 & = OB^2 + OC^2 - 2 OB \cdot OC \cos \angle BOC
\end{align*}

代入得
\begin{align*}
    \cos \angle OA & = \dfrac{2OB \cdot OC \cos \angle BOC - 2OA^2}{2AB \cdot AC} \\
    & = \dfrac{\cos \angle BOC - \frac{OA}{OB} \frac{OA}{OC}}{\frac{AB}{OB} \frac{AC}{OC}} \\
    & = \dfrac{\cos \angle BOC - \cos \angle AOB \cos \angle AOC}{\sin \angle AOB \sin \angle AOC}
\end{align*}

这样,我们就得到了如下定理.

\begin{thm}[三面角第一余弦定理]
    设$OA, OB, OC$为空间中共顶点的三条射线,则所成三面角的二面角为
    \begin{align}
        \cos \angle OA & = \dfrac{\cos \angle BOC - \cos \angle AOB \cos \angle AOC}{\sin \angle AOB \sin \angle AOC} \label{1} \\
        \cos \angle OB & = \dfrac{\cos \angle AOC - \cos \angle AOB \cos \angle BOC}{\sin \angle AOB \sin \angle BOC} \label{2} \\
        \cos \angle OC & = \dfrac{\cos \angle AOB - \cos \angle BOC \cos \angle AOC}{\sin \angle BOC \sin \angle AOC} \label{3}
    \end{align}
\end{thm}

反过来,我们也可以由二面角的余弦值求出平面角的余弦值.

\begin{thm}[三面角第二余弦定理]
    设$OA, OB, OC$为空间中共顶点的三条射线,则所成三面角的平面角为
    \begin{align*}
        \cos \angle AOB & = \dfrac{\cos \angle OC + \cos \angle OA \cos \angle OB}{\sin \angle OA \sin \angle OB} \\
        \cos \angle BOC & = \dfrac{\cos \angle OA + \cos \angle OB \cos \angle OC}{\sin \angle OB \sin \angle OC} \\
        \cos \angle AOC & = \dfrac{\cos \angle OB + \cos \angle OC \cos \angle OA}{\sin \angle OC \sin \angle OA}
    \end{align*}
\end{thm}

\begin{proof}
    将 \eqref{3} 改写为
    \begin{equation}
        \cos \angle AOB  = \cos \angle OC \sin \angle BOC \sin \angle AOC + \cos \angle BOC \cos \angle AOC \label{4}
    \end{equation}

    另外,由 \eqref{1}, \eqref{2} 可得
    \begin{align*}
        \sin \angle BOC & = \dfrac{\sin \angle OA}{\sin \angle AOB} \\
        \sin \angle AOC & = \dfrac{\sin \angle OB}{\sin \angle AOB}
    \end{align*}

    代入 \eqref{4} 可得
    \begin{align*}
        \cos \angle AOB & = \cos \angle OC \cdot \dfrac{\sin \angle OA}{\sin \angle AOB} \cdot
            \dfrac{\sin \angle OB}{\sin \angle AOB} + \cos \angle BOC \cos \angle AOC \\
        & = \dfrac{\cos \angle OC \sin \angle OA \sin \angle OB + \cos \angle OA \cos \angle OB \sin^2 \angle AOB}
            {\sin^2 \angle AOB}
    \end{align*}

    注意到$\sin^2 \angle AOB = 1 - \cos^2 \angle AOB$,代入整理得
    \[ \cos \angle AOB = \dfrac{\cos \angle OC + \cos \angle OA \cos \angle OB}{\sin \angle OA \sin \angle OB} \]

    类似可得剩余两式.
\end{proof}

\end{document}